\documentclass[11pt]{article}
\usepackage[a4paper,margin=24mm]{geometry}
\usepackage{amsmath,amssymb,amsthm,booktabs,hyperref}
\newcommand{\E}{\mathbb E}\newcommand{\Pp}{\mathbb P}
\newcommand{\CE}{C_{\rm E}}
\newtheorem{theorem}{Theorem}\newtheorem{corollary}{Corollary}
\title{Continuous low-noise confinement without a probability grid}
\author{Bartosz Ko\l{}odziejek \and Kacper Rodziewicz}\date{28 October 2026}
\begin{document}\maketitle
\noindent\textit{AI-assisted research note.} The argument and exact scalar
checks were developed with AI assistance and reviewed separately. The
result below is conditional on its stated analytic and entry premises;
it is not a proof of the full sharp iid Berry--Esseen conjecture.

\section{Conditional statement and acquired application}
Let $n\ge42$ be an actual integer. Let $(P,H)$ be a fixed-target maximizing
pair with nonnegative excess $\mathfrak a$ over $c=\CE$, where $X\sim P$
is centered and has variance one. Retain the closed-support contact theorem
and the true predecessor bound $C_{n-1}\le c$. The actual raw threshold
$H$ can have either sign. Assume an actual central gap $(-A,B)$ with
$A,B>0$, and write
\[
 p=\Pp(X>0),\quad q=1-p,\quad m=\E X_+=\E X_-,\quad
 a=m/q,\quad b=m/p,\quad \ell=a+b,
\]
\[
 X=\ell(Y-p),\qquad Y=\mathcal B+U,\qquad
 \mathcal B=\mathbf1_{\{X>0\}},\qquad \E(U\mid\mathcal B)=0.
\]
Define $U_i$ as the conditional value of $U$ on $\{\mathcal B=i\}$.
Assume the already acquired bounds
\begin{equation}\label{ln:entry}
 p\in[2/5,3/5],\qquad |U|<3/5,\qquad
 d=1-(A+B)/\ell<9/125,
 \qquad \lambda=n\E U^2\le3/40.
\end{equation}
The gap implies $A\le a$, $B\le b$, and hence $d\ge0$.

\begin{theorem}\label{ln:theorem}
Under these hypotheses, the entire conditional supports satisfy
$|U|<1/2$. The conclusion is uniform in all raw $H$ and all $n\ge42$.
It includes the closed noise endpoint $\lambda=3/40$ and zero noise.
\end{theorem}

The probability interval in this abstract lemma is $[0.4,0.6]$.
The currently acquired entry supplies the incoming strip and inward bounds
only on the union of panels 40--52, namely $[0.40,0.53]$. The application
below is restricted to that smaller interval. In particular this note
does not acquire the strip on $(0.53,0.6]$ or treat the larger-noise region.

\section{One common moment enclosure}
Put $v=pq$, $s=\E U^2=\lambda/n$, and
$\theta=\E(X-\E[X\mid\mathcal B])^2$. Conditional centering and variance
give the exact identities
\begin{equation}\label{ln:identities}
 \ell^2=(v+s)^{-1},\qquad
 \theta=\frac{s}{v+s}=\frac{\lambda}{nv+\lambda},\qquad
 \ell^2=\frac{1-\theta}{v}.
\end{equation}
Throughout \eqref{ln:entry}, $6/25\le v\le1/4$ and
$\max(p,q)\le3/5$. Use the rational constants
\[
 \Theta=\frac{37}{5000},\qquad E=\frac{249}{125},\qquad
 r=\frac{221}{200}.
\]
Then
\[
 \theta\le\frac{5}{677}<\Theta,\qquad
 \ell^2\ge4(1-\Theta)>E^2,
\]
with exact differences
$\Theta-5/677=49/3385000$ and $4(1-\Theta)-E^2=73/31250$.

Because $\mathcal B$ is the actual sign of $X$, the absolute cubic
expansion has no sign approximation:
\begin{equation}\label{ln:cubic}
 \rho:=\E|X|^3=\ell^3\left\{v(1-2v)
       +3\E[|\mathcal B-p|U^2]
       +\E[\operatorname{sgn}(\mathcal B-p)U^3]\right\}.
\end{equation}
The linear term vanishes by conditional centering. The incoming strip,
used before the conclusion, gives $|U|^3\le(3/5)U^2$ and
$|\mathcal B-p|\le3/5$. Also $\ell^3\le v^{-3/2}$. Therefore
\[
 \rho\le\rho_B(v)+\frac{12/5}{v^{3/2}}\frac\lambda n,
 \qquad \rho_B(v)=\frac{1-2v}{\sqrt v}.
\]
The function $\rho_B$ decreases on $(0,1/4]$, since its derivative is
$-\tfrac12v^{-3/2}-v^{-1/2}<0$. The exact squared comparison
\[
 24^2(6/25)^3-(12/5)^2=\frac{34416}{15625}>0
\]
gives $\rho\le\rho_B(6/25)+24\lambda/n<r$. For the last step,
\[
 r-\frac3{70}=\frac{1487}{1400}>0,\qquad
 \left(r-\frac3{70}\right)^2\frac6{25}
       -\left(\frac{13}{25}\right)^2=\frac{8707}{24500000}>0.
\]

If $\theta_\pm$ are the two nonnegative conditional variance contributions
in physical coordinates, their sum is $\theta$. Expanding the signed
second moment gives
\[
 \beta:=\E[X|X|]=(1-\theta)(1-2p)+\theta_+-\theta_-.
\]
Consequently, for the actual $p$, including $p>1/2$,
\begin{equation}\label{ln:beta}
 1-2p-2(1-p)\Theta\le\beta\le1-2p+2p\Theta.
\end{equation}
All these bounds concern one common marginal. No assertion about
independent conditional laws or a moment relaxation is used.

\section{All-threshold contact and exterior rays}
For $w=n-1$, $s_n=\sqrt{n/w}$, $F_w(t)=\Phi(t/\sqrt w)$ and
$e_n=n(s_n-1)-1/2$, fixed-target contact and the predecessor bound give,
at every closed-support point $t$,
\begin{equation}\label{ln:contact}
 0\le\mathcal H_n(H,t)
       +c\{-|t|^3+\tfrac32\rho t^2+3\beta t+\rho e_n\}
       -n\mathfrak a,
\end{equation}
where
\[
 \frac{\mathcal H_n(H,t)}{n^{3/2}}
 =F_w(H-t)-F_n(H)+tF_n'(H)-\frac{t^2-1}{2}F_n''(H).
\]
The term $-n\mathfrak a$ is nonpositive and may be dropped. This is a
fixed-target statement; no equality between $c$ and the actual score is
assumed.

Expand the first term spatially through degree four. Expand the other
three terms in variance through orders two, one and one, using
$\partial_w F_w=F_w''/2$. At $t=\pm u$, cancellations give
\[
 \frac{\mathcal H_n(H,\pm u)}{s_n^3}
 =\pm A_u(1-z^2)\phi(z)
       +\frac{H_4(u)}{24\sqrt w}\phi'''(z)+R,
\]
\[
 z=H/\sqrt w,\quad A_u=(u^3-3u)/6,
 \qquad H_4(u)=u^4-6u^2+3.
\]
The absolute bounds for the four remainders are
\[
 \frac{u^5\|\phi^{(4)}\|_\infty}{120w},\quad
 \frac{\|\phi^{(5)}\|_\infty}{48w^{3/2}},\quad
 \frac{u\|\phi^{(4)}\|_\infty}{8w},\quad
 \frac{|u^2-1|\|\phi^{(5)}\|_\infty}{16w^{3/2}}.
\]
Gaussian Fourier inversion, integrating absolute powers of frequency,
gives $\|\phi^{(4)}\|_\infty\le3\phi(0)<6/5$ and
$\|\phi^{(5)}\|_\infty\le8/\pi<51/20$.
For $u\ge1$ their sum is at most
\[
 \frac{u^5/100+3u/20}{w}
      +\frac{51(3u^2-2)}{960w^{3/2}}.
\]
Thus the remainder extends below the earlier $u\ge5/2$ presentation.
For $u>\sqrt3$, $A_u>0$, and the elementary bounds
$\|(1-z^2)\phi(z)\|_\infty=\phi(0)$ and
$\|\phi'''\|_\infty\le2/\pi<2\phi(0)$ control both signs.

Set $S=81/80$, $c_0=819/2000<c$ and $e=8/875$.
For $n\ge42$, $s_n<S$, $\sqrt w>6$, and $e_n<e$.
For completeness,
$ne_n=s_n^2(2s_n+1)/(2(s_n+1)^2)$ is increasing for $s_n\ge1$;
its value at $S$ is $793881/2073680<48/125$.
Define
\begin{align}
 E_u&=u^5/4100+3u/820+17(3u^2-2)/78720,\notag\\
 Q_{r,b}(u)&=-u^3+\tfrac32ru^2+3bu+er,\label{ln:Q}\\
 P_{r,b}(u)&=\tfrac25S^3\{A_u+|H_4(u)|/72\}
                  +S^3 E_u+c_0Q_{r,b}(u).\label{ln:P}
\end{align}
If the directional signed moment is at most $b$, $Q<0$ licenses replacing
the positive multiplier $c$ by $c_0$. Then $P<0$ contradicts
\eqref{ln:contact} at that boundary for every raw $H$.

To continue along a physical ray, hold $r,b$ fixed. Direct differentiation
of the exact heat expression gives, with $z=H/\sqrt n$ and
$y=(H-t)/\sqrt n$,
\[
 \partial_t\mathcal H_n=n\phi(z)+Ht\phi(z)-ns_n\phi(s_ny).
\]
Taylor's theorem and $\|\phi''\|_\infty=\phi(0)$ bound the spatial
remainder by $\phi(0)t^2/2$. Since
$\partial_s[s\phi(sy)]=-\phi''(sy)$, the variance contribution is at
most $\phi(0)n(s_n-1)$. Hence
\[
 |\partial_t\mathcal H_n|<\tfrac25(t^2/2+1/2+e).
\]
Put $L_{r,b}(u)=\tfrac25(u^2/2+1/2+e)+c_0Q'_{r,b}(u)$.
If $Q,Q',Q'',L,L',L''$ are negative at a boundary, the signs persist
along the whole ray: $Q'''=-6$ and $L''=2/5-6c_0<0$.
Thus the contact upper bound strictly decreases there. Boundary atoms
and the entire closed exterior ray are excluded.

\section{A continuous one-parameter certificate}
For the negative ray use $z=p$ and the signed moment $-\beta$; for the
positive ray use $z=1-p$ and $\beta$. In both cases $z\in[2/5,3/5]$,
and \eqref{ln:beta} gives the directional cap
\[
 b_z=2z-1+2(1-z)\Theta.
\]
This folds the actual probability, not the endpoints of a probability
panel. It does not reflect the chosen discrepancy or change its law.
Take the boundary $u=E(z+1/2)$. Thus
\[
 u\in[u_0,u_1]=[2241/1250,2739/1250],\qquad
 b_z=\alpha u+\delta,
 \quad\alpha=\frac{4963}{4980},\quad\delta=-\frac{9889}{5000}.
\]
Here $u_0^2>3$, so $H_4'(u)=4u(u^2-3)>0$;
$H_4(u_1)<-2.75$ proves $H_4<0$ on this whole interval.
Let $\mathcal P(u)=P_{r,\alpha u+\delta}(u)$.
Replacing $|H_4|$ by $-H_4$ in \eqref{ln:P} gives an explicit quintic
$\mathcal P(u)=\sum_{j=0}^5 a_j u^j$, with
\begin{align*}
 a_5&=S^3/4100,& a_4&=-\tfrac25S^3/72,\\
 a_3&=\tfrac25S^3/6-c_0,&
 a_2&=\tfrac25S^3/12+51S^3/78720+c_0(3r/2+3\alpha),\\
 a_1&=-\tfrac25S^3/2+3S^3/820+3c_0\delta,&
 a_0&=-\tfrac25S^3/24-34S^3/78720+c_0er.
\end{align*}
The exact rational endpoint inequalities are
\[
 \mathcal P'''(u_0)<-2.24,\qquad \mathcal P'''(u_1)<-2.27,
 \qquad \mathcal P'(u_0)>0.915,\quad \mathcal P'(u_1)>0.745.
\]
Indeed $\mathcal P'''$ is a convex quadratic, with second derivative
$120a_5>0$. Its negative endpoint values make it negative throughout.
Therefore $\mathcal P'$ is concave and exceeds the chord between its
positive endpoint values. It follows that $\mathcal P$ increases and
\begin{equation}\label{ln:margin}
 \mathcal P(u)\le\mathcal P(u_1)
 =-\frac{4560620039489147528734839}{25625000000000000000000000}
 <-0.1779.
\end{equation}
One stronger-looking shortcut is false:
$\mathcal P''(u_0)>0.0228$. It is not used.

The multiplier and ray signs require a separate calculation.
The boundary family
\[
 \mathcal Q(u)=Q_{r,\alpha u+\delta}(u)
 =-u^3+(3r/2+3\alpha)u^2+3\delta u+er
\]
has total derivatives $\mathcal Q''(u_0)<-1.46$ and
$\mathcal Q'(u_1)=179223/6250000>0.0286$.
Since $\mathcal Q''$ decreases, $\mathcal Q'$ decreases and stays
positive. Hence $\mathcal Q$ increases, and
$\mathcal Q(u)\le\mathcal Q(u_1)<-1.198$.

For each physical ray the directional cap is fixed. Its derivatives,
evaluated across the boundary family, are instead
\begin{align*}
 q_1(u)&=-3u^2+(3r+3\alpha)u+3\delta,&
 q_2(u)&=-6u+3r,\\
 l(u)&=\tfrac25(u^2/2+1/2+e)+c_0q_1(u),&
 l_1(u)&=\tfrac25u+c_0q_2(u),\qquad l_2=\tfrac25-6c_0.
\end{align*}
They are not the total derivatives of $\mathcal Q$.
The following rational upper bounds certify their signs on the interval:
\begin{center}
\begin{tabular}{lrr}\toprule
family & value at $u_0$ & total derivative at $u_0$\\\midrule
$q_1$ & $<-4.272$ & $<-4.452$\\
$q_2$ & $<-7.441$ & $=-6$\\
$l$   & $<-0.903$ & $<-1.105$\\
$l_1$ & $<-2.330$ & $=-2.057$\\\bottomrule
\end{tabular}
\end{center}
The two quadratic families have negative constant second derivatives
$q_1''=-6$ and $l''=-2.057$; the others are affine decreasing.
Thus all four stay negative, as does $l_2=-2.057$.
Together with $\mathcal Q<0$, these are the six required signs at every
boundary. Holding its actual directional cap fixed, the preceding ray
argument excludes the whole exterior ray beyond each boundary.

\section{Completion, application and remaining premises}
For each actual $p$ the two excluded rays start at
\[
 u_-=E(p+1/2),\qquad u_+=E(3/2-p).
\]
Because $\ell>E$, the remaining support satisfies
$-1/2<Y=X/\ell+p<3/2$ strictly. The already proved actual gap supplies
\[
 U_0\le(a-A)/\ell\le d<9/125,
 \qquad U_1\ge-(b-B)/\ell\ge-d>-9/125.
\]
For $\mathcal B=0$ the new lower bound is $U_0>-1/2$;
for $\mathcal B=1$ the new upper bound is $U_1<1/2$.
This proves Theorem~\ref{ln:theorem}. If $\lambda=0$, then $\E U^2=0$
and $U=0$ directly. There is no division by noise or limiting argument.

\begin{corollary}\label{ln:application}
Retain the accepted fixed-target entry and the ordered elementary
moment/gap contraction on its 30 probability panels. On its acquired
subrange $p\in[0.40,0.53]$, all actual $n\ge42$ with
$\lambda\in[0,3/40]$ satisfy $|U|<1/2$ for every raw $H$.
The previously acquired inward bound $d<0.072<1/4$ is retained.
\end{corollary}
\begin{proof}
That antecedent provides the strip $|U|<0.6$ on the 13 closed panels
40--52 and $d\le0.0719539227<0.072$.
Their closed union is $[0.40,0.53]$. Apply the theorem at the actual $p$.
\end{proof}

The new continuous calculation uses exact rational arithmetic and 28
domain/endpoint inequalities, with no new probability, phase or noise
grid. Six intentionally false conditions are rejected, including the
incorrect second-derivative shortcut and substitution of a total
derivative for a fixed-cap ray derivative. The earlier 13-row calculation
and its fresh bounded antecedent replay are preserved as additional
evidence; those heat rows are not premises of this continuous proof.

The mathematical contact theorem, predecessor bound and incoming
moment/gap acquisition remain separate premises. This result neither
certifies a finite base nor supplies strict count exclusion. It does not
remove any existing certificate or assert full-noise confinement on all
13 panels. The already separate full-noise result on $[0.41,0.42]$ is
unchanged.
\end{document}
